\section{多刚体系统树形拓扑描述}

\subsection{动力学模型拓扑图}

以四足机器人为例，其多刚体系统动力学模型如图~\ref{fig:10.1}~所示。

\begin{figure}[htbp]
  \centering
  \caption{四足机器人动力学模型}
  \label{fig:10.1}
\end{figure}

动力学模型各元件连接关系可用图~\ref{fig:10.2}~所示的动力学模型拓扑图表示。

\begin{figure}[htbp]
  \centering
  \caption{四足机器人动力学模型拓扑图}
  \label{fig:10.2}
\end{figure}

\subsection{派生树拓扑描述与元件标号}

通过切断系统中的闭环，可将四足机器人多刚体系统动力学模型处理为派生树拓扑结构；编号时确保任意体元件的标号大于其内接体的标号；非切断铰元件编号与其外接体元件编号相同，切断铰元件编号与其对应的闭环相同，如图~\ref{fig:10.3}~所示。

\begin{figure}[htbp]
  \centering
  \caption{四足机器人动力学模型机身与右前腿元件标号}
  \label{fig:10.3}
\end{figure}

各元件内接体编号以及前序元件集合如表~\ref{tab:10.1}~所示。

\begin{table}[htbp]
  \centering
  \caption{各元件内接体编号及其前序元件集合}
  \label{tab:10.1}
  \begin{tabular}{ccc}
    \hline
    $k$ & $\vec{p}(k)$ & $\vec{p}[k]$ \\
    \hline
    $1$  & $0$  & $\{1\}$ \\
    $\vdots$ & $\vdots$ & $\vdots$ \\
    $11$ & $1$  & $\{1,11\}$ \\
    $12$ & $11$ & $\{1,11,12\}$ \\
    $13$ & $12$ & $\{1,11,12,13\}$ \\
    $\vdots$ & $\vdots$ & $\vdots$ \\
    $20$ & $12$ & $\{1,11,12,20\}$ \\
    $21$ & $20$ & $\{1,11,12,20,21\}$ \\
    \hline
  \end{tabular}
\end{table}

任意体元件 $k$ 的前序元件集合满足如下递推规律
\begin{equation}\label{eq:10.1}
  \vec{p}[k] = \begin{cases}
    \emptyset, & k = 0, \\
    \{\vec{p}[p(k)],\, k\}, & k \neq 0.
  \end{cases}
\end{equation}

四足机器人完整的派生树及元件编号如图~\ref{fig:10.4}~所示。

\begin{figure}[htbp]
  \centering
  \caption{四足机器人树形拓扑图}
  \label{fig:10.4}
\end{figure}

四足机器人动力学模型中每一个闭环对应着一个切断铰，与这些切断铰关联的元件序号如表~\ref{tab:10.2}~所示。

\begin{table}[htbp]
  \centering
  \caption{四足机器人闭环切断铰关联的体元件}
  \label{tab:10.2}
  \begin{tabular}{cccc}
    \hline
    $j$ & $b(j)$ & $m(j)$ & $r(j)$ \\
    \hline
    $1$ & $4$  & $15$ & $3$ \\
    $2$ & $7$  & $17$ & $6$ \\
    $3$ & $10$ & $19$ & $9$ \\
    $4$ & $13$ & $31$ & $12$ \\
    \hline
  \end{tabular}
\end{table}

\subsection{系统广义坐标与约束方程}

系统广义坐标 $\boldsymbol{q}$ 由系统派生树中非切断铰元件的相对运动广义坐标组成，并满足由切断铰元件引起的闭环约束方程 $\boldsymbol{0} = \boldsymbol{c}(\boldsymbol{q})$，可分别表示为
\begin{equation}\label{eq:10.2}
  \boldsymbol{q} = \begin{bmatrix} \boldsymbol{q}_{J_1} \\ \boldsymbol{q}_{J_2} \\ \vdots \\ \boldsymbol{q}_{J_n} \end{bmatrix}, \qquad
  \boldsymbol{c} = \begin{bmatrix} \boldsymbol{c}_{L_1} \\ \boldsymbol{c}_{L_2} \\ \vdots \\ \boldsymbol{c}_{L_m} \end{bmatrix},
\end{equation}
式中 $n$ 为非切断铰元件的个数（等于系统包含的体元件个数）；$m$ 为切断铰元件的个数（等于系统中闭环的个数）。对于前述四足机器人多刚体系统动力学模型：$n = 21$，$m = 4$。

\section{相邻体元件运动学量递推关系}

\subsection{相邻体元件位置级运动学量递推关系}

以 $O_i$ 为原点的体元件连体坐标系简称为 $i$ 系。将大地元件连体坐标系，即 $0$ 系，视为全局惯性坐标系，则 $i\;(i=1,2,\ldots,n)$ 系在全局惯性坐标系中的位置和方位可通过递推的方式表示为
\begin{equation}\label{eq:10.3}
  \begin{cases}
    \boldsymbol{r}_{O_0 O_i \mid 0} = \boldsymbol{r}_{O_0 O_{p(i)} \mid 0} + \boldsymbol{A}_{0,p(i)} \boldsymbol{r}_{O_{p(i)} O_i \mid p(i)}, \\
    \boldsymbol{A}_{0,i} = \boldsymbol{A}_{0,p(i)} \boldsymbol{A}_{p(i),i},
  \end{cases}
\end{equation}
式中 $\boldsymbol{r}_{P_1 P_2 \mid i}$ 表示从 $P_1$ 点到 $P_2$ 的位置矢量在 $i$ 系中的投影，$\boldsymbol{A}_{i,j}$ 表示从 $j$ 系到 $i$ 系的坐标变换矩阵，$\boldsymbol{r}_{O_{p(i)} O_i \mid p(i)}$ 又可表示为
\begin{equation}\label{eq:10.4}
  \boldsymbol{r}_{O_{p(i)} O_i \mid p(i)} = \boldsymbol{r}_{O_{p(i)} R_i \mid p(i)} + \boldsymbol{r}_{R_i O_i \mid p(i)}.
\end{equation}
对于刚体元件 $i$，$\boldsymbol{r}_{O_{p(i)} R_i \mid p(i)}$ 为定常矩阵，其时间导数为零。以上等式中 $\boldsymbol{r}_{R_i O_i \mid p(i)}$ 和 $\boldsymbol{A}_{p(i),i}$ 均可表示为铰元件 $J_i$ 广义坐标 $\boldsymbol{q}_{J_i}$ 的表达式。

\subsection{相邻体元件速度级运动学量递推关系}

将 $j$ 系相对于 $i$ 系的相对平动速度和相对转动速度在 $j$ 系中的六维投影列阵记为 $\boldsymbol{v}_{i,j \mid j}$，即
\begin{equation}\label{eq:10.5}
  \boldsymbol{v}_{i,j \mid j} \triangleq \begin{bmatrix} \boldsymbol{A}_{i,j}^{\mathrm{T}} \dot{\boldsymbol{r}}_{O_i O_j \mid i} \\ \boldsymbol{\omega}_{i,j \mid j} \end{bmatrix},
\end{equation}
式中 $\boldsymbol{\omega}_{i,j \mid j}$ 表示 $j$ 系相对于 $i$ 系的相对角速度矢量在 $j$ 系中的投影；$\dot{\boldsymbol{r}}$ 表示 $\boldsymbol{r}$ 对时间的一阶导数。相邻元件速度运动学量递推表达式可记为
\begin{equation}\label{eq:10.6}
  \boldsymbol{v}_{0,i \mid i} = \boldsymbol{C}_{i,p(i)} \boldsymbol{v}_{0,p(i) \mid p(i)} + \boldsymbol{\Phi}_{J_i} \dot{\boldsymbol{q}}_{J_i},
\end{equation}
相关证明见~11.11.1~小节，式中
\begin{equation}\label{eq:10.7}
  \boldsymbol{C}_{i,p(i)} \triangleq \begin{bmatrix}
    \boldsymbol{A}_{p(i),i}^{\mathrm{T}} & \boldsymbol{A}_{p(i),i}^{\mathrm{T}} \tilde{\boldsymbol{r}}_{O_{p(i)} O_i \mid p(i)}^{\mathrm{T}} \\
    \boldsymbol{O}_3 & \boldsymbol{A}_{p(i),i}^{\mathrm{T}}
  \end{bmatrix}.
\end{equation}
对于刚体元件，$\boldsymbol{C}_{i,p(i)}$ 仅与 $p(i)$ 系和 $i$ 系的相对位置和相对方位有关，可记为铰元件 $J_i$ 广义坐标 $\boldsymbol{q}_{J_i}$ 的表达式，如图~\ref{fig:10.5}~所示；$\boldsymbol{\Phi}_{J_i}$ 也为与 $\boldsymbol{q}_{J_i}$ 相关的表达式。

\begin{figure}[htbp]
  \centering
  \caption{牵连速度系数矩阵}
  \label{fig:10.5}
\end{figure}

\subsection{相邻体元件加速度运动学量递推关系}

将 $j$ 系相对于 $i$ 系的相对平动加速度和相对转动加速度在 $j$ 系中的六维投影列阵记为 $\boldsymbol{a}_{i,j \mid j}$，即
\begin{equation}\label{eq:10.8}
  \boldsymbol{a}_{i,j \mid j} \triangleq \begin{bmatrix} \boldsymbol{A}_{i,j}^{\mathrm{T}} \ddot{\boldsymbol{r}}_{O_i O_j \mid i} \\ \dot{\boldsymbol{\omega}}_{i,j \mid j} \end{bmatrix},
\end{equation}
相邻元件加速度运动学量递推表达式可记为
\begin{equation}\label{eq:10.9}
  \boldsymbol{a}_{0,i \mid i} = \boldsymbol{C}_{i,p(i)} \boldsymbol{a}_{0,p(i) \mid p(i)} + \boldsymbol{\Phi}_{J_i} \ddot{\boldsymbol{q}}_{J_i} + \boldsymbol{\eta}_i,
\end{equation}
相关证明见~11.11.2~小节，式中
\begin{equation}\label{eq:10.10}
  \boldsymbol{\eta}_i = \boldsymbol{\xi}_{J_i} + \begin{bmatrix}
    \boldsymbol{A}_{p(i),i}^{\mathrm{T}} \left( \tilde{\boldsymbol{\omega}}_{0,p(i) \mid p(i)} \tilde{\boldsymbol{\omega}}_{0,p(i) \mid p(i)} \boldsymbol{r}_{O_{p(i)} O_i \mid p(i)} + 2 \tilde{\boldsymbol{\omega}}_{0,p(i) \mid p(i)} \dot{\boldsymbol{r}}_{O_{p(i)} O_i \mid p(i)} \right) \\[4pt]
    \boldsymbol{A}_{p(i),i}^{\mathrm{T}} \tilde{\boldsymbol{\omega}}_{0,p(i) \mid p(i)} \boldsymbol{\omega}_{p(i),i \mid p(i)}
  \end{bmatrix},
\end{equation}
式中各类铰元件相对加速度非齐次项 $\boldsymbol{\xi}_{J_i}$ 的具体表达式见~\ref{sec:10.3}~节。

\section{主铰元件相对运动学方程}\label{sec:10.3}

\subsection{虚拟铰元件}

对于虚拟铰元件 $i$，位置级相对运动学方程可表示为
\begin{equation}\label{eq:10.11}
  \begin{cases}
    \boldsymbol{r}_{R_i O_i \mid p(i)} = \boldsymbol{l}_{J_i}, \\
    \boldsymbol{A}_{p(i),i} = \boldsymbol{A}_Z(\theta_{J_i,1}) \boldsymbol{A}_Y(\theta_{J_i,2}) \boldsymbol{A}_X(\theta_{J_i,3}),
  \end{cases}
\end{equation}
式中 $\begin{bmatrix} x_{J_i} & y_{J_i} & z_{J_i} \end{bmatrix}^{\mathrm{T}} =: \boldsymbol{l}_{J_i}$ 和 $\begin{bmatrix} \theta_{J_i,1} & \theta_{J_i,2} & \theta_{J_i,3} \end{bmatrix}^{\mathrm{T}} =: \boldsymbol{\theta}_{J_i}$ 分别为描述虚拟铰元件 $i$ 平动和转动的广义坐标，绕坐标轴的单位旋转矩阵表达式分别为
\begin{equation}\label{eq:10.11b}
  \boldsymbol{A}_X(\theta) = \begin{bmatrix} 1 & 0 & 0 \\ 0 & \cos\theta & -\sin\theta \\ 0 & \sin\theta & \cos\theta \end{bmatrix}, \quad
  \boldsymbol{A}_Y(\theta) = \begin{bmatrix} \cos\theta & 0 & \sin\theta \\ 0 & 1 & 0 \\ -\sin\theta & 0 & \cos\theta \end{bmatrix}, \quad
  \boldsymbol{A}_Z(\theta) = \begin{bmatrix} \cos\theta & -\sin\theta & 0 \\ \sin\theta & \cos\theta & 0 \\ 0 & 0 & 1 \end{bmatrix}.
\end{equation}

速度级相对运动学方程可表示为
\begin{equation}\label{eq:10.12}
  \begin{cases}
    \dot{\boldsymbol{r}}_{R_i O_i \mid p(i)} = \dot{\boldsymbol{l}}_{J_i}, \\
    \boldsymbol{\omega}_{p(i),i \mid i} = \boldsymbol{H}_{B321}(\boldsymbol{\theta}_{J_i}) \dot{\boldsymbol{\theta}}_{J_i},
  \end{cases}
\end{equation}
由此可得
\begin{equation}\label{eq:10.13}
  \boldsymbol{v}_{p(i),i \mid i} = \begin{bmatrix} \boldsymbol{A}_{p(i),i}^{\mathrm{T}} \dot{\boldsymbol{r}}_{O_{p(i)} O_i \mid p(i)} \\ \boldsymbol{\omega}_{p(i),i \mid i} \end{bmatrix}
  = \begin{bmatrix} \boldsymbol{A}_{p(i),i}^{\mathrm{T}} \dot{\boldsymbol{r}}_{R_i O_i \mid p(i)} \\ \boldsymbol{\omega}_{p(i),i \mid i} \end{bmatrix}
  = \begin{bmatrix} \boldsymbol{A}_{p(i),i}^{\mathrm{T}} & \boldsymbol{O}_3 \\ \boldsymbol{O}_3 & \boldsymbol{H}_{B321}(\boldsymbol{\theta}_{J_i}) \end{bmatrix}
  \begin{bmatrix} \dot{\boldsymbol{l}}_{J_i} \\ \dot{\boldsymbol{\theta}}_{J_i} \end{bmatrix}
  =: \boldsymbol{\Phi}_{J_i} \dot{\boldsymbol{q}}_{J_i}.
\end{equation}

加速度相对运动学方程可表示为
\begin{equation}\label{eq:10.14}
  \begin{cases}
    \ddot{\boldsymbol{r}}_{R_i O_i \mid p(i)} = \ddot{\boldsymbol{l}}_{J_i}, \\
    \dot{\boldsymbol{\omega}}_{p(i),i \mid i} = \boldsymbol{H}_{B321}(\boldsymbol{\theta}_{J_i}) \ddot{\boldsymbol{\theta}}_{J_i} + \dot{\boldsymbol{H}}_{B321}(\boldsymbol{\theta}_{J_i}, \dot{\boldsymbol{\theta}}_{J_i}) \dot{\boldsymbol{\theta}}_{J_i},
  \end{cases}
\end{equation}
由此可得
\begin{equation}\label{eq:10.15}
  \boldsymbol{a}_{p(i),i \mid i} = \begin{bmatrix} \boldsymbol{A}_{p(i),i}^{\mathrm{T}} \ddot{\boldsymbol{r}}_{O_{p(i)} O_i \mid p(i)} \\ \dot{\boldsymbol{\omega}}_{p(i),i \mid i} \end{bmatrix}
  = \begin{bmatrix} \boldsymbol{A}_{p(i),i}^{\mathrm{T}} & \boldsymbol{O}_3 \\ \boldsymbol{O}_3 & \boldsymbol{H}_{B321}(\boldsymbol{\theta}_{J_i}) \end{bmatrix}
  \begin{bmatrix} \ddot{\boldsymbol{l}}_{J_i} \\ \ddot{\boldsymbol{\theta}}_{J_i} \end{bmatrix}
  + \begin{bmatrix} \boldsymbol{0}_3 \\ \dot{\boldsymbol{H}}_{B321}(\boldsymbol{\theta}_{J_i}, \dot{\boldsymbol{\theta}}_{J_i}) \dot{\boldsymbol{\theta}}_{J_i} \end{bmatrix}
  =: \boldsymbol{\Phi}_{J_i} \ddot{\boldsymbol{q}}_{J_i} + \boldsymbol{\xi}_{J_i}.
\end{equation}

上述等式中
\begin{equation}\label{eq:10.15b}
  \boldsymbol{H}_{B321}(\boldsymbol{\theta}) = \begin{bmatrix}
    -s_2 & 0 & 1 \\
    c_2 s_3 & c_3 & 0 \\
    c_2 c_3 & -s_3 & 0
  \end{bmatrix}, \quad
  \dot{\boldsymbol{H}}_{B321} = \begin{bmatrix}
    -c_2 \dot{\theta}_2 & 0 & 0 \\
    -s_2 s_3 \dot{\theta}_2 + c_2 c_3 \dot{\theta}_3 & -s_3 \dot{\theta}_3 & 0 \\
    -s_2 c_3 \dot{\theta}_2 - c_2 s_3 \dot{\theta}_3 & -c_3 \dot{\theta}_3 & 0
  \end{bmatrix},
\end{equation}
式中 $c_2 = \cos(\theta_2)$，$s_2 = \sin(\theta_2)$，$c_3 = \cos(\theta_3)$，$s_3 = \sin(\theta_3)$。

\subsection{转动铰元件}

对于转动铰元件 $i$，位置级相对运动学方程可表示为
\begin{equation}\label{eq:10.16}
  \begin{cases}
    \boldsymbol{r}_{R_i O_i \mid p(i)} = \boldsymbol{0}_3, \\
    \boldsymbol{A}_{p(i),i} = \boldsymbol{I}_3 \cos\theta_{J_i} + \tilde{\boldsymbol{n}}_{J_i} \sin\theta_{J_i} + \boldsymbol{n}_{J_i} \boldsymbol{n}_{J_i}^{\mathrm{T}} (1 - \cos\theta_{J_i}) \\
    \phantom{\boldsymbol{A}_{p(i),i} =} = \boldsymbol{I}_3 + \tilde{\boldsymbol{n}}_{J_i} \sin\theta_{J_i} + \tilde{\boldsymbol{n}}_{J_i} \tilde{\boldsymbol{n}}_{J_i} (1 - \cos\theta_{J_i}),
  \end{cases}
\end{equation}
式中 $\boldsymbol{n}_{J_i}$ 表示转动铰元件 $i$ 沿转轴方向的单位矢量在 $i$ 系和 $p(i)$ 系中的投影列阵。

速度级相对运动学方程可表示为
\begin{equation}\label{eq:10.17}
  \begin{cases}
    \dot{\boldsymbol{r}}_{R_i O_i \mid p(i)} = \boldsymbol{0}_3, \\
    \boldsymbol{\omega}_{p(i),i \mid i} = \boldsymbol{n}_{J_i} \dot{\theta}_{J_i},
  \end{cases}
\end{equation}
由此可得
\begin{equation}\label{eq:10.18}
  \boldsymbol{v}_{p(i),i \mid i} = \begin{bmatrix} \boldsymbol{A}_{p(i),i}^{\mathrm{T}} \dot{\boldsymbol{r}}_{O_{p(i)} O_i \mid p(i)} \\ \boldsymbol{\omega}_{p(i),i \mid i} \end{bmatrix}
  = \begin{bmatrix} \boldsymbol{0}_3 \\ \boldsymbol{n}_{J_i} \end{bmatrix} \dot{\theta}_{J_i}
  =: \boldsymbol{\Phi}_{J_i} \dot{\boldsymbol{q}}_{J_i}.
\end{equation}

加速度相对运动学方程可表示为
\begin{equation}\label{eq:10.19}
  \begin{cases}
    \ddot{\boldsymbol{r}}_{R_i O_i \mid p(i)} = \boldsymbol{0}_3, \\
    \dot{\boldsymbol{\omega}}_{p(i),i \mid i} = \boldsymbol{n}_{J_i} \ddot{\theta}_{J_i},
  \end{cases}
\end{equation}
由此可得
\begin{equation}\label{eq:10.20}
  \boldsymbol{a}_{p(i),i \mid i} = \begin{bmatrix} \boldsymbol{A}_{p(i),i}^{\mathrm{T}} \ddot{\boldsymbol{r}}_{O_{p(i)} O_i \mid p(i)} \\ \dot{\boldsymbol{\omega}}_{p(i),i \mid i} \end{bmatrix}
  = \begin{bmatrix} \boldsymbol{0}_3 \\ \boldsymbol{n}_{J_i} \end{bmatrix} \ddot{\theta}_{J_i}
  =: \boldsymbol{\Phi}_{J_i} \ddot{\boldsymbol{q}}_{J_i} + \boldsymbol{\xi}_{J_i}.
\end{equation}

\subsection{棱柱铰元件}

对于棱柱铰元件 $i$，位置级相对运动学方程可表示为
\begin{equation}\label{eq:10.21}
  \begin{cases}
    \boldsymbol{r}_{R_i O_i \mid p(i)} = \boldsymbol{n}_{J_i} l_{J_i}, \\
    \boldsymbol{A}_{p(i),i} = \boldsymbol{I}_3,
  \end{cases}
\end{equation}
式中 $\boldsymbol{n}_{J_i}$ 表示棱柱铰元件 $i$ 沿滑移轴方向的单位矢量在 $i$ 系和 $p(i)$ 系中的投影列阵。

速度级相对运动学方程可表示为
\begin{equation}\label{eq:10.22}
  \begin{cases}
    \dot{\boldsymbol{r}}_{R_i O_i \mid p(i)} = \boldsymbol{n}_{J_i} \dot{l}_{J_i}, \\
    \boldsymbol{\omega}_{p(i),i \mid i} = \boldsymbol{0},
  \end{cases}
\end{equation}
由此可得
\begin{equation}\label{eq:10.23}
  \boldsymbol{v}_{p(i),i \mid i} = \begin{bmatrix} \boldsymbol{A}_{p(i),i}^{\mathrm{T}} \dot{\boldsymbol{r}}_{O_{p(i)} O_i \mid p(i)} \\ \boldsymbol{\omega}_{p(i),i \mid i} \end{bmatrix}
  = \begin{bmatrix} \boldsymbol{n}_{J_i} \\ \boldsymbol{0}_3 \end{bmatrix} \dot{l}_{J_i}
  =: \boldsymbol{\Phi}_{J_i} \dot{\boldsymbol{q}}_{J_i}.
\end{equation}

加速度相对运动学方程可表示为
\begin{equation}\label{eq:10.24}
  \begin{cases}
    \ddot{\boldsymbol{r}}_{R_i O_i \mid p(i)} = \boldsymbol{n}_{J_i} \ddot{l}_{J_i}, \\
    \dot{\boldsymbol{\omega}}_{p(i),i \mid i} = \boldsymbol{0}_3,
  \end{cases}
\end{equation}
由此可得
\begin{equation}\label{eq:10.25}
  \boldsymbol{a}_{p(i),i \mid i} = \begin{bmatrix} \boldsymbol{A}_{p(i),i}^{\mathrm{T}} \ddot{\boldsymbol{r}}_{O_{p(i)} O_i \mid p(i)} \\ \dot{\boldsymbol{\omega}}_{p(i),i \mid i} \end{bmatrix}
  = \begin{bmatrix} \boldsymbol{n}_{J_i} \\ \boldsymbol{0}_3 \end{bmatrix} \ddot{l}_{J_i}
  =: \boldsymbol{\Phi}_{J_i} \ddot{\boldsymbol{q}}_{J_i} + \boldsymbol{\xi}_{J_i}.
\end{equation}

\subsection{固结铰元件}

对于固结铰元件 $i$，位置级相对运动学方程可表示为
\begin{equation}\label{eq:10.26}
  \begin{cases}
    \boldsymbol{r}_{R_i O_i \mid p(i)} = \boldsymbol{0}_3, \\
    \boldsymbol{A}_{p(i),i} = \boldsymbol{I}_3.
  \end{cases}
\end{equation}

速度级相对运动学方程可表示为
\begin{equation}\label{eq:10.27}
  \begin{cases}
    \dot{\boldsymbol{r}}_{R_i O_i \mid p(i)} = \boldsymbol{0}_3, \\
    \boldsymbol{\omega}_{p(i),i \mid i} = \boldsymbol{0}_3,
  \end{cases}
\end{equation}
由此可得
\begin{equation}\label{eq:10.28}
  \boldsymbol{v}_{p(i),i \mid i} = \begin{bmatrix} \boldsymbol{A}_{p(i),i}^{\mathrm{T}} \dot{\boldsymbol{r}}_{O_{p(i)} O_i \mid p(i)} \\ \boldsymbol{\omega}_{p(i),i \mid i} \end{bmatrix}
  = \begin{bmatrix} \boldsymbol{0}_3 \\ \boldsymbol{0}_3 \end{bmatrix}.
\end{equation}

加速度相对运动学方程可表示为
\begin{equation}\label{eq:10.29}
  \begin{cases}
    \ddot{\boldsymbol{r}}_{R_i O_i \mid p(i)} = \boldsymbol{0}_3, \\
    \dot{\boldsymbol{\omega}}_{p(i),i \mid i} = \boldsymbol{0}_3,
  \end{cases}
\end{equation}
由此可得
\begin{equation}\label{eq:10.30}
  \boldsymbol{a}_{p(i),i \mid i} = \begin{bmatrix} \boldsymbol{A}_{p(i),i}^{\mathrm{T}} \ddot{\boldsymbol{r}}_{O_{p(i)} O_i \mid p(i)} \\ \dot{\boldsymbol{\omega}}_{p(i),i \mid i} \end{bmatrix}
  = \begin{bmatrix} \boldsymbol{0}_3 \\ \boldsymbol{0}_3 \end{bmatrix}.
\end{equation}

\subsection{球铰元件}

对于球铰元件 $i$，位置级相对运动学方程可表示为
\begin{equation}\label{eq:10.31}
  \begin{cases}
    \boldsymbol{r}_{R_i O_i \mid p(i)} = \boldsymbol{0}, \\
    \boldsymbol{A}_{p(i),i} = \boldsymbol{A}_Z(\theta_{J_i,1}) \boldsymbol{A}_Y(\theta_{J_i,2}) \boldsymbol{A}_X(\theta_{J_i,3}),
  \end{cases}
\end{equation}
式中 $\begin{bmatrix} \theta_{J_i,1} & \theta_{J_i,2} & \theta_{J_i,3} \end{bmatrix}^{\mathrm{T}} =: \boldsymbol{\theta}_{J_i}$ 为描述球铰元件 $i$ 相对转动的广义坐标，绕坐标轴的单位旋转矩阵表达式参见~\ref{sec:10.3}~节。

速度级相对运动学方程可表示为
\begin{equation}\label{eq:10.32}
  \begin{cases}
    \dot{\boldsymbol{r}}_{R_i O_i \mid p(i)} = \boldsymbol{0}, \\
    \boldsymbol{\omega}_{p(i),i \mid i} = \boldsymbol{H}_{B321}(\boldsymbol{\theta}_{J_i}) \dot{\boldsymbol{\theta}}_{J_i},
  \end{cases}
\end{equation}
由此可得
\begin{equation}\label{eq:10.33}
  \boldsymbol{v}_{p(i),i \mid i} = \begin{bmatrix} \boldsymbol{A}_{p(i),i}^{\mathrm{T}} \dot{\boldsymbol{r}}_{O_{p(i)} O_i \mid p(i)} \\ \boldsymbol{\omega}_{p(i),i \mid i} \end{bmatrix}
  = \begin{bmatrix} \boldsymbol{O}_3 \\ \boldsymbol{H}_{B321}(\boldsymbol{\theta}_{J_i}) \end{bmatrix} \dot{\boldsymbol{\theta}}_{J_i}
  =: \boldsymbol{\Phi}_{J_i} \dot{\boldsymbol{q}}_{J_i}.
\end{equation}

加速度相对运动学方程可表示为
\begin{equation}\label{eq:10.34}
  \begin{cases}
    \ddot{\boldsymbol{r}}_{R_i O_i \mid p(i)} = \boldsymbol{0}, \\
    \dot{\boldsymbol{\omega}}_{p(i),i \mid i} = \boldsymbol{H}_{B321}(\boldsymbol{\theta}_{J_i}) \ddot{\boldsymbol{\theta}}_{J_i} + \dot{\boldsymbol{H}}_{B321}(\boldsymbol{\theta}_{J_i}, \dot{\boldsymbol{\theta}}_{J_i}) \dot{\boldsymbol{\theta}}_{J_i},
  \end{cases}
\end{equation}
由此可得
\begin{equation}\label{eq:10.35}
  \boldsymbol{a}_{p(i),i \mid i} = \begin{bmatrix} \boldsymbol{A}_{p(i),i}^{\mathrm{T}} \ddot{\boldsymbol{r}}_{O_{p(i)} O_i \mid p(i)} \\ \dot{\boldsymbol{\omega}}_{p(i),i \mid i} \end{bmatrix}
  = \begin{bmatrix} \boldsymbol{O}_3 \\ \boldsymbol{H}_{B321}(\boldsymbol{\theta}_{J_i}) \end{bmatrix} \ddot{\boldsymbol{\theta}}_{J_i}
  + \begin{bmatrix} \boldsymbol{0}_3 \\ \dot{\boldsymbol{H}}_{B321}(\boldsymbol{\theta}_{J_i}, \dot{\boldsymbol{\theta}}_{J_i}) \dot{\boldsymbol{\theta}}_{J_i} \end{bmatrix}
  =: \boldsymbol{\Phi}_{J_i} \ddot{\boldsymbol{q}}_{J_i} + \boldsymbol{\xi}_{J_i}.
\end{equation}

上述等式中 $\boldsymbol{H}_{B321}$ 和 $\dot{\boldsymbol{H}}_{B321}$ 的具体表达式详见~\ref{sec:10.3}~节。

\subsection{万向节}

万向节（universal joint）允许两根相互垂直轴的转动，其位置级、速度级和加速度级相对运动学方程可仿照转动铰元件的推导方法，结合两个转动自由度进行建立。

\subsection{圆柱铰}

圆柱铰（cylindrical joint）同时允许沿轴线的平动和绕轴线的转动，其相对运动学方程可结合棱柱铰与转动铰的运动学方程进行建立。

\subsection{螺旋铰}

螺旋铰（screw joint）的平动与转动满足固定螺距约束关系，其相对运动学方程可在圆柱铰的基础上增加螺距约束进行建立。

\section{体元件速度关于系统广义速度的线性表达式}

对于任意体元件 $i$，其绝对速度列阵 $\boldsymbol{v}_{0,i \mid i}$ 可记为如下形式
\begin{equation}\label{eq:10.37}
  \boldsymbol{v}_{0,i \mid i} = \sum_{k \in \vec{p}[i]} \boldsymbol{C}_{i,k} \boldsymbol{\Phi}_{J_k} \dot{\boldsymbol{q}}_{J_k},
\end{equation}
式中牵连速度矩阵
\begin{equation}\label{eq:10.38}
  \boldsymbol{C}_{i,k} = \begin{bmatrix}
    \boldsymbol{A}_{k,i}^{\mathrm{T}} & \boldsymbol{A}_{k,i}^{\mathrm{T}} \tilde{\boldsymbol{r}}_{O_k O_i \mid k}^{\mathrm{T}} \\
    \boldsymbol{O}_3 & \boldsymbol{A}_{k,i}^{\mathrm{T}}
  \end{bmatrix}, \quad k \in \vec{p}[i],
\end{equation}
证明过程见~11.11.3~小节，示意图如图~\ref{fig:10.6}~所示。当 $k = p(i)$ 时，有
\begin{equation}\label{eq:10.39}
  \boldsymbol{C}_{i,p(i)} = \begin{bmatrix}
    \boldsymbol{A}_{p(i),i}^{\mathrm{T}} & \boldsymbol{A}_{p(i),i}^{\mathrm{T}} \tilde{\boldsymbol{r}}_{O_{p(i)} O_i \mid p(i)}^{\mathrm{T}} \\
    \boldsymbol{O}_3 & \boldsymbol{A}_{p(i),i}^{\mathrm{T}}
  \end{bmatrix},
\end{equation}
该表达式与式~(\ref{eq:10.7})~相同；当 $k = i$ 时，有
\begin{equation}\label{eq:10.40}
  \boldsymbol{C}_{i,i} = \begin{bmatrix}
    \boldsymbol{A}_{i,i}^{\mathrm{T}} & \boldsymbol{A}_{i,i}^{\mathrm{T}} \tilde{\boldsymbol{r}}_{O_i O_i \mid i}^{\mathrm{T}} \\
    \boldsymbol{O}_3 & \boldsymbol{A}_{i,i}^{\mathrm{T}}
  \end{bmatrix}
  = \begin{bmatrix} \boldsymbol{I}_3 & \boldsymbol{O}_3 \\ \boldsymbol{O}_3 & \boldsymbol{I}_3 \end{bmatrix}
  = \boldsymbol{I}_6.
\end{equation}

对于任意元件 $i$，其牵连速度系数矩阵~(\ref{eq:10.38})~满足如下递推规律（运算量较大）
\begin{equation}\label{eq:10.41}
  \boldsymbol{C}_{i,k} = \begin{cases}
    \boldsymbol{I}_6, & k = i, \\
    \boldsymbol{C}_{i,p(i)} \boldsymbol{C}_{p(i),k}, & k \in \vec{p}[\vec{p}[i]].
  \end{cases}
\end{equation}

\begin{figure}[htbp]
  \centering
  \caption{体元件前序元件牵连速度系数矩阵}
  \label{fig:10.6}
\end{figure}

\section{闭环切断铰元件相对运动学量的计算}

系统闭环切断铰元件相对运动学量并未作为系统广义坐标的一部分，需借助当前时刻的系统广义坐标来确定，计算过程中有时也要综合考虑该切断铰上一时刻的相对运动学量，以防计算结果发生突变。

\subsection{转动铰元件}

当 $L_j$ 为转动铰时，动系到基系的坐标转换矩阵可解读为
\begin{equation}\label{eq:10.42}
  \boldsymbol{A}_{b(j),m(j)} = \boldsymbol{I}_3 \cos\theta_{L_j} + \tilde{\boldsymbol{n}}_{L_j} \sin\theta_{L_j} + \boldsymbol{n}_{L_j} \boldsymbol{n}_{L_j}^{\mathrm{T}} (1 - \cos\theta_{L_j})
  =: \begin{bmatrix} a_{11} & a_{12} & a_{13} \\ a_{21} & a_{22} & a_{23} \\ a_{31} & a_{32} & a_{33} \end{bmatrix}_{b(j),m(j)}.
\end{equation}
由此可得
\begin{equation}\label{eq:10.43}
  \mathrm{tr}\,\boldsymbol{A}_{b(j),m(j)} = 1 + 2\cos\theta_{L_j}
\end{equation}
和
\begin{equation}\label{eq:10.44}
  2 \boldsymbol{n}_{L_j} \sin\theta_{L_j} = \begin{bmatrix} a_{32} - a_{23} \\ a_{13} - a_{31} \\ a_{21} - a_{12} \end{bmatrix}_{b(j),m(j)}.
\end{equation}
由以上两式进一步可得
\begin{equation}\label{eq:10.45}
  \cos\theta_{L_j} = \frac{\mathrm{tr}\,\boldsymbol{A}_{b(j),m(j)} - 1}{2} =: X
\end{equation}
和
\begin{equation}\label{eq:10.46}
  \sin\theta_{L_j} = \frac{1}{2\boldsymbol{n}_{L_j}^{\mathrm{T}}} \begin{bmatrix} a_{32} - a_{23} \\ a_{13} - a_{31} \\ a_{21} - a_{12} \end{bmatrix}_{b(j),m(j)} =: Y.
\end{equation}
在 $(-\pi,\,\pi]$ 区间内可得
\begin{equation}\label{eq:10.47}
  \hat{\theta}_{L_j} = \mathrm{atan2}(Y,\, X).
\end{equation}

将上一时刻 $L_j$ 的转角记为
\begin{equation}\label{eq:10.48}
  \theta_{L_j}^- = \hat{\theta}_{L_j}^- + 2\pi N_{L_j}^-,
\end{equation}
式中 $\hat{\theta}_{L_j}^- \in (-\pi,\,\pi]$，$N_{L_j}^- \in \mathbb{N}$。则当前时刻运动学量满足
\begin{equation}\label{eq:10.49}
  \hat{\theta}_{L_j}^+ = \hat{\theta}_{L_j},
\end{equation}
\begin{equation}\label{eq:10.50}
  N_{L_j}^+ = \begin{cases}
    N_{L_j}^- + 1, & \hat{\theta}_{L_j}^- > \pi/2 \;\text{and}\; \hat{\theta}_{L_j} < -\pi/2, \\
    N_{L_j}^- - 1, & \hat{\theta}_{L_j}^- < -\pi/2 \;\text{and}\; \hat{\theta}_{L_j} > \pi/2,
  \end{cases}
\end{equation}
\begin{equation}\label{eq:10.51}
  \theta_{L_j} = \hat{\theta}_{L_j}^+ + 2\pi N_{L_j}^+.
\end{equation}

\subsection{棱柱铰元件}

当切断铰 $L_j$ 为棱柱铰时，可由位置级约束方程直接求解滑移位移 $l_{L_j}$，再通过对时间求导获得速度和加速度。

\subsection{固结铰元件}

当切断铰 $L_j$ 为固结铰时，切断铰两侧体元件的相对位置和相对方位均为固定值，无需额外广义坐标。

\subsection{球铰元件}

当切断铰 $L_j$ 为球铰时，可由坐标转换矩阵 $\boldsymbol{A}_{b(j),m(j)}$ 反解出三个转动广义坐标。

\subsection{万向节}

当切断铰 $L_j$ 为万向节时，可由坐标转换矩阵反解出两个转动广义坐标。

\section{闭环切断铰元件约束方程}

\subsection{转动铰元件}

当闭环 $L_j$ 切断铰为转动铰元件时，其位置级约束方程可表示为
\begin{equation}\label{eq:10.52}
  \boldsymbol{0} = \boldsymbol{c}_{L_j} = \begin{bmatrix}
    \boldsymbol{r}_{O_{r(j)} M_j \mid r(j)} - \boldsymbol{r}_{O_{r(j)} B_j \mid r(j)} \\[4pt]
    \boldsymbol{T}_{L_j}^{\mathrm{T}} \boldsymbol{A}_{r(j),b(j)}^{\mathrm{T}} \boldsymbol{A}_{r(j),m(j)} \boldsymbol{n}_{L_j}
  \end{bmatrix}
  =: \begin{bmatrix} \boldsymbol{c}_{L_j,U} \\ \boldsymbol{c}_{L_j,D} \end{bmatrix},
\end{equation}
式中
\begin{align}
  \boldsymbol{r}_{O_{r(j)} M_j \mid r(j)} &= \boldsymbol{r}_{O_{r(j)} O_{m(j)} \mid r(j)} + \boldsymbol{A}_{r(j),m(j)} \boldsymbol{r}_{O_{m(j)} M_j \mid m(j)}, \label{eq:10.53} \\
  \boldsymbol{r}_{O_{r(j)} B_j \mid r(j)} &= \boldsymbol{r}_{O_{r(j)} O_{b(j)} \mid r(j)} + \boldsymbol{A}_{r(j),b(j)} \boldsymbol{r}_{O_{b(j)} B_j \mid b(j)}, \label{eq:10.54}
\end{align}
三维列阵 $\boldsymbol{n}_{L_j}$ 为在 $b(j)$ 系和 $m(j)$ 系中表示的切断铰转轴方向的单位矢量。$\boldsymbol{T}_{L_j} = \begin{bmatrix} \boldsymbol{\tau}_{L_j,1} & \boldsymbol{\tau}_{L_j,2} \end{bmatrix}$，$\boldsymbol{\tau}_{L_j,1}$ 和 $\boldsymbol{\tau}_{L_j,2}$ 为建模时定义的一对模为 $1$ 的定常三维列阵；$\boldsymbol{\tau}_{L_j,1}$、$\boldsymbol{\tau}_{L_j,2}$ 与 $\boldsymbol{n}_{L_j}$ 两两正交。

位置级约束方程中 $\boldsymbol{c}_{L_j,U}$ 关于时间的一阶导数满足
\begin{align}
  \dot{\boldsymbol{c}}_{L_j,U} &= \dot{\boldsymbol{r}}_{O_{r(j)} M_j \mid r(j)} - \dot{\boldsymbol{r}}_{O_{r(j)} B_j \mid r(j)} \label{eq:10.55} \\
  &= \frac{\mathrm{d}}{\mathrm{d}t}\left( \boldsymbol{r}_{O_{r(j)} O_{m(j)} \mid r(j)} + \boldsymbol{A}_{r(j),m(j)} \boldsymbol{r}_{O_{m(j)} M_j \mid m(j)} \right) - \frac{\mathrm{d}}{\mathrm{d}t}\left( \boldsymbol{r}_{O_{r(j)} O_{b(j)} \mid r(j)} + \boldsymbol{A}_{r(j),b(j)} \boldsymbol{r}_{O_{b(j)} B_j \mid b(j)} \right) \notag \\
  &= \dot{\boldsymbol{r}}_{O_{r(j)} O_{m(j)} \mid r(j)} + \boldsymbol{A}_{r(j),m(j)} \tilde{\boldsymbol{\omega}}_{r(j),m(j) \mid m(j)} \boldsymbol{r}_{O_{m(j)} M_j \mid m(j)} \notag \\
  &\quad - \dot{\boldsymbol{r}}_{O_{r(j)} O_{b(j)} \mid r(j)} - \boldsymbol{A}_{r(j),b(j)} \tilde{\boldsymbol{\omega}}_{r(j),b(j) \mid b(j)} \boldsymbol{r}_{O_{b(j)} B_j \mid b(j)} \notag \\
  &= \boldsymbol{A}_{r(j),m(j)} \boldsymbol{A}_{r(j),m(j)}^{\mathrm{T}} \dot{\boldsymbol{r}}_{O_{r(j)} O_{m(j)} \mid r(j)} + \boldsymbol{A}_{r(j),m(j)} \tilde{\boldsymbol{r}}_{O_{m(j)} M_j \mid m(j)}^{\mathrm{T}} \boldsymbol{\omega}_{r(j),m(j) \mid m(j)} \notag \\
  &\quad - \boldsymbol{A}_{r(j),b(j)} \boldsymbol{A}_{r(j),b(j)}^{\mathrm{T}} \dot{\boldsymbol{r}}_{O_{r(j)} O_{b(j)} \mid r(j)} - \boldsymbol{A}_{r(j),b(j)} \tilde{\boldsymbol{r}}_{O_{b(j)} B_j \mid b(j)}^{\mathrm{T}} \boldsymbol{\omega}_{r(j),b(j) \mid b(j)} \notag \\
  &= \begin{bmatrix} \boldsymbol{A}_{r(j),m(j)} & \boldsymbol{A}_{r(j),m(j)} \tilde{\boldsymbol{r}}_{O_{m(j)} M_j \mid m(j)}^{\mathrm{T}} \end{bmatrix}
  \begin{bmatrix} \boldsymbol{A}_{r(j),m(j)}^{\mathrm{T}} \dot{\boldsymbol{r}}_{O_{r(j)} O_{m(j)} \mid r(j)} \\ \boldsymbol{\omega}_{r(j),m(j) \mid m(j)} \end{bmatrix} \notag \\
  &\quad - \begin{bmatrix} \boldsymbol{A}_{r(j),b(j)} & \boldsymbol{A}_{r(j),b(j)} \tilde{\boldsymbol{r}}_{O_{b(j)} B_j \mid b(j)}^{\mathrm{T}} \end{bmatrix}
  \begin{bmatrix} \boldsymbol{A}_{r(j),b(j)}^{\mathrm{T}} \dot{\boldsymbol{r}}_{O_{r(j)} O_{b(j)} \mid r(j)} \\ \boldsymbol{\omega}_{r(j),b(j) \mid b(j)} \end{bmatrix} \notag \\
  &=: \begin{bmatrix} \boldsymbol{A}_{r(j),m(j)} & \boldsymbol{A}_{r(j),m(j)} \tilde{\boldsymbol{r}}_{O_{m(j)} M_j \mid m(j)}^{\mathrm{T}} \end{bmatrix} \boldsymbol{v}_{r(j),m(j) \mid m(j)} \notag \\
  &\quad - \begin{bmatrix} \boldsymbol{A}_{r(j),b(j)} & \boldsymbol{A}_{r(j),b(j)} \tilde{\boldsymbol{r}}_{O_{b(j)} B_j \mid b(j)}^{\mathrm{T}} \end{bmatrix} \boldsymbol{v}_{r(j),b(j) \mid b(j)}. \label{eq:10.56}
\end{align}

位置级约束方程中 $\boldsymbol{c}_{L_j,D}$ 关于时间的一阶导数满足
\begin{align}
  \dot{\boldsymbol{c}}_{L_j,D} &= \frac{\mathrm{d}}{\mathrm{d}t}\left( \boldsymbol{T}_{L_j}^{\mathrm{T}} \boldsymbol{A}_{r(j),b(j)}^{\mathrm{T}} \boldsymbol{A}_{r(j),m(j)} \boldsymbol{n}_{L_j} \right) \label{eq:10.57} \\
  &= \boldsymbol{T}_{L_j}^{\mathrm{T}} \dot{\boldsymbol{A}}_{r(j),b(j)}^{\mathrm{T}} \boldsymbol{A}_{r(j),m(j)} \boldsymbol{n}_{L_j} + \boldsymbol{T}_{L_j}^{\mathrm{T}} \boldsymbol{A}_{r(j),b(j)}^{\mathrm{T}} \dot{\boldsymbol{A}}_{r(j),m(j)} \boldsymbol{n}_{L_j} \notag \\
  &= \boldsymbol{T}_{L_j}^{\mathrm{T}} \boldsymbol{A}_{r(j),b(j)}^{\mathrm{T}} \tilde{\boldsymbol{\omega}}_{r(j),b(j) \mid r(j)}^{\mathrm{T}} \boldsymbol{A}_{r(j),m(j)} \boldsymbol{n}_{L_j} - \boldsymbol{T}_{L_j}^{\mathrm{T}} \boldsymbol{A}_{r(j),b(j)}^{\mathrm{T}} \tilde{\boldsymbol{\omega}}_{r(j),m(j) \mid r(j)}^{\mathrm{T}} \boldsymbol{A}_{r(j),m(j)} \boldsymbol{n}_{L_j} \notag \\
  &= \boldsymbol{T}_{L_j}^{\mathrm{T}} \boldsymbol{A}_{r(j),b(j)}^{\mathrm{T}} \boldsymbol{A}_{r(j),m(j)} \tilde{\boldsymbol{n}}_{L_j} \boldsymbol{\omega}_{r(j),b(j) \mid m(j)} - \boldsymbol{T}_{L_j}^{\mathrm{T}} \boldsymbol{A}_{r(j),b(j)}^{\mathrm{T}} \boldsymbol{A}_{r(j),m(j)} \tilde{\boldsymbol{n}}_{L_j} \boldsymbol{\omega}_{r(j),m(j) \mid m(j)} \notag \\
  &= \boldsymbol{T}_{L_j}^{\mathrm{T}} \boldsymbol{A}_{b(j),m(j)} \tilde{\boldsymbol{n}}_{L_j} \boldsymbol{\omega}_{r(j),b(j) \mid m(j)} - \boldsymbol{T}_{L_j}^{\mathrm{T}} \boldsymbol{A}_{b(j),m(j)} \tilde{\boldsymbol{n}}_{L_j} \boldsymbol{\omega}_{r(j),m(j) \mid m(j)} \notag \\
  &= \boldsymbol{T}_{L_j}^{\mathrm{T}} \boldsymbol{A}_{b(j),m(j)} \tilde{\boldsymbol{n}}_{L_j} \boldsymbol{A}_{b(j),m(j)}^{\mathrm{T}} \boldsymbol{\omega}_{r(j),b(j) \mid b(j)} - \boldsymbol{T}_{L_j}^{\mathrm{T}} \boldsymbol{A}_{b(j),m(j)} \tilde{\boldsymbol{n}}_{L_j} \boldsymbol{\omega}_{r(j),m(j) \mid m(j)} \notag \\
  &= \begin{bmatrix} \boldsymbol{O}_{2\times 3} & \boldsymbol{T}_{L_j}^{\mathrm{T}} \boldsymbol{A}_{b(j),m(j)} \tilde{\boldsymbol{n}}_{L_j}^{\mathrm{T}} \end{bmatrix}
  \begin{bmatrix} \boldsymbol{A}_{r(j),m(j)}^{\mathrm{T}} \dot{\boldsymbol{r}}_{O_{r(j)} O_{m(j)} \mid r(j)} \\ \boldsymbol{\omega}_{r(j),m(j) \mid m(j)} \end{bmatrix} \notag \\
  &\quad - \begin{bmatrix} \boldsymbol{O}_{2\times 3} & \boldsymbol{T}_{L_j}^{\mathrm{T}} \boldsymbol{A}_{b(j),m(j)} \tilde{\boldsymbol{n}}_{L_j}^{\mathrm{T}} \boldsymbol{A}_{b(j),m(j)}^{\mathrm{T}} \end{bmatrix}
  \begin{bmatrix} \boldsymbol{A}_{r(j),b(j)}^{\mathrm{T}} \dot{\boldsymbol{r}}_{O_{r(j)} O_{b(j)} \mid r(j)} \\ \boldsymbol{\omega}_{r(j),b(j) \mid b(j)} \end{bmatrix} \notag \\
  &= \begin{bmatrix} \boldsymbol{O}_{2\times 3} & \boldsymbol{T}_{L_j}^{\mathrm{T}} \boldsymbol{A}_{b(j),m(j)} \tilde{\boldsymbol{n}}_{L_j}^{\mathrm{T}} \end{bmatrix} \boldsymbol{v}_{r(j),m(j) \mid m(j)} \notag \\
  &\quad - \begin{bmatrix} \boldsymbol{O}_{2\times 3} & \boldsymbol{T}_{L_j}^{\mathrm{T}} \boldsymbol{A}_{b(j),m(j)} \tilde{\boldsymbol{n}}_{L_j}^{\mathrm{T}} \boldsymbol{A}_{b(j),m(j)}^{\mathrm{T}} \end{bmatrix} \boldsymbol{v}_{r(j),b(j) \mid b(j)}. \label{eq:10.57b}
\end{align}

当闭环 $L_j$ 切断铰为转动铰元件时，则有
\begin{equation}\label{eq:10.58}
  \dot{\boldsymbol{c}}_{L_j} = \begin{bmatrix} \dot{\boldsymbol{c}}_{L_j,U} \\ \dot{\boldsymbol{c}}_{L_j,D} \end{bmatrix}
  = \begin{bmatrix}
    \boldsymbol{A}_{r(j),m(j)} & \boldsymbol{A}_{r(j),m(j)} \tilde{\boldsymbol{r}}_{O_{m(j)} M_j \mid m(j)}^{\mathrm{T}} \\
    \boldsymbol{O}_{2\times 3} & \boldsymbol{T}_{L_j}^{\mathrm{T}} \boldsymbol{A}_{b(j),m(j)} \tilde{\boldsymbol{n}}_{L_j}^{\mathrm{T}}
  \end{bmatrix} \boldsymbol{v}_{r(j),m(j) \mid m(j)}
  - \begin{bmatrix}
    \boldsymbol{A}_{r(j),b(j)} & \boldsymbol{A}_{r(j),b(j)} \tilde{\boldsymbol{r}}_{O_{b(j)} B_j \mid b(j)}^{\mathrm{T}} \\
    \boldsymbol{O}_{2\times 3} & \boldsymbol{T}_{L_j}^{\mathrm{T}} \boldsymbol{A}_{b(j),m(j)} \tilde{\boldsymbol{n}}_{L_j}^{\mathrm{T}} \boldsymbol{A}_{b(j),m(j)}^{\mathrm{T}}
  \end{bmatrix} \boldsymbol{v}_{r(j),b(j) \mid b(j)}
\end{equation}
\begin{equation}\label{eq:10.58b}
  =: \boldsymbol{\Psi}_{m(j)} \boldsymbol{v}_{r(j),m(j) \mid m(j)} - \boldsymbol{\Psi}_{b(j)} \boldsymbol{v}_{r(j),b(j) \mid b(j)}.
\end{equation}

\subsection{棱柱铰元件}

当切断铰 $L_j$ 为棱柱铰时，约束方程仅含位置约束（转轴方向不变），其速度级方程可仿照转动铰的推导方法建立。

\subsection{固结铰元件}

当切断铰 $L_j$ 为固结铰时，约束方程包含完整的位置约束和方向约束。

\subsection{球铰元件}

当切断铰 $L_j$ 为球铰时，约束方程仅含位置约束（无方向约束），其速度级方程形式相应简化。

\subsection{万向节}

当切断铰 $L_j$ 为万向节时，约束方程包含位置约束和一个方向约束。

\section{速度级约束方程关于系统广义速度的线性表达式}

对于任意闭环切断铰 $L_j$，有
\begin{align}
  \dot{\boldsymbol{c}}_{L_j} &= \boldsymbol{\Psi}_{m(j)} \boldsymbol{v}_{r(j),m(j) \mid m(j)} - \boldsymbol{\Psi}_{b(j)} \boldsymbol{v}_{r(j),b(j) \mid b(j)} \label{eq:10.59} \\
  &= \boldsymbol{\Psi}_{m(j)} \sum_{\substack{k \in \vec{p}[m(j)] \\ k > r(j)}} \boldsymbol{C}_{m(j),k} \boldsymbol{\Phi}_{J_k} \dot{\boldsymbol{q}}_{J_k} - \boldsymbol{\Psi}_{b(j)} \sum_{\substack{k \in \vec{p}[b(j)] \\ k > r(j)}} \boldsymbol{C}_{b(j),k} \boldsymbol{\Phi}_{J_k} \dot{\boldsymbol{q}}_{J_k} \notag \\
  &= \sum_{\substack{k \in \vec{p}[m(j)] \\ k > r(j)}} \boldsymbol{\Psi}_{m(j)} \boldsymbol{C}_{m(j),k} \boldsymbol{\Phi}_{J_k} \dot{\boldsymbol{q}}_{J_k} - \sum_{\substack{k \in \vec{p}[b(j)] \\ k > r(j)}} \boldsymbol{\Psi}_{b(j)} \boldsymbol{C}_{b(j),k} \boldsymbol{\Phi}_{J_k} \dot{\boldsymbol{q}}_{J_k} \notag \\
  &=: \sum_{\substack{k \in \vec{p}[m(j)] \\ k > r(j)}} \boldsymbol{c}_{L_j,q_{J_k}} \dot{\boldsymbol{q}}_{J_k} + \sum_{\substack{k \in \vec{p}[b(j)] \\ k > r(j)}} \boldsymbol{c}_{L_j,q_{J_k}} \dot{\boldsymbol{q}}_{J_k} \notag \\
  &=: \boldsymbol{c}_{L_j,q} \dot{\boldsymbol{q}} + \boldsymbol{c}_{L_j,t}.
\end{align}

上述推导过程用到了表达式
\begin{equation}\label{eq:10.60}
  \{k \mid k \in \vec{p}[m(j)],\, k > r(j)\} \cap \{k \mid k \in \vec{p}[b(j)],\, k > r(j)\} = \emptyset
\end{equation}
和
\begin{equation}\label{eq:10.61}
  \frac{\partial \dot{\boldsymbol{c}}_{L_j}}{\partial \dot{\boldsymbol{q}}_{J_k}} \equiv \frac{\partial \boldsymbol{c}_{L_j}}{\partial \boldsymbol{q}_{J_k}} =: \boldsymbol{c}_{L_j,q_{J_k}}.
\end{equation}

\section{闭环违约修正}

导致闭环约束方程出现违约的因素包括初始装配偏差、浮点数舍入误差、微分方程数值解法截断误差以及累积误差等。

\subsection{位置级违约修正}

在任意给定时刻 $t$，针对一组给定的系统广义坐标 $\boldsymbol{q}$，当位置级约束方程违约量的按模最大值超过给定的阈值时，即当
\begin{equation}\label{eq:10.62}
  \|\boldsymbol{c}(\boldsymbol{q}, t)\|_\infty > \varepsilon
\end{equation}
时，需修正系统广义坐标 $\boldsymbol{q}$，即寻找系统广义坐标的修正量 $\Delta\boldsymbol{q}$，使得修正后的系统广义坐标 $\boldsymbol{q} + \Delta\boldsymbol{q}$ 对应的违约量 $\boldsymbol{c}(\boldsymbol{q} + \Delta\boldsymbol{q}, t)$ 满足阈值要求，即
\begin{equation}\label{eq:10.63}
  \|\boldsymbol{c}(\boldsymbol{q} + \Delta\boldsymbol{q}, t)\|_\infty \le \varepsilon.
\end{equation}

上述非线性不等式方程不易求解，在满足该非线性不等式方程的前提下，不妨令
\begin{equation}\label{eq:10.64}
  \boldsymbol{c}(\boldsymbol{q} + \Delta\boldsymbol{q}, t) = \boldsymbol{0}.
\end{equation}
上式为关于 $\Delta\boldsymbol{q}$ 的非线性方程组，仍不易直接求解，可进一步作如下线性化处理，即
\begin{equation}\label{eq:10.65}
  \boldsymbol{0} = \boldsymbol{c}(\boldsymbol{q} + \Delta\boldsymbol{q}, t) = \boldsymbol{c}(\boldsymbol{q}, t) + \boldsymbol{c}_q(\boldsymbol{q}, t) \Delta\boldsymbol{q} + o(\Delta\boldsymbol{q}).
\end{equation}
忽略 $\Delta\boldsymbol{q}$ 的二阶及二阶以上小量，则有
\begin{equation}\label{eq:10.66}
  \boldsymbol{0} = \boldsymbol{c}(\boldsymbol{q}, t) + \boldsymbol{c}_q(\boldsymbol{q}, t) \Delta\boldsymbol{q}.
\end{equation}
一般情况下，上述关于 $\Delta\boldsymbol{q}$ 的线性方程组有无穷多组解，在这些解中寻求二范数最小的一组解，求解过程不妨简记为如下表达式
\begin{equation}\label{eq:10.67}
  \Delta\boldsymbol{q} = -\boldsymbol{c}_q^+ \boldsymbol{c}.
\end{equation}
根据求得的系统广义坐标修正量 $\Delta\boldsymbol{q}$ 更新系统广义坐标，即
\begin{equation}\label{eq:10.68}
  \boldsymbol{q} \leftarrow \boldsymbol{q} + \Delta\boldsymbol{q}.
\end{equation}
针对更新后的系统广义坐标，再次判定位置级约束方程违约量的按模最大值是否超过给定的阈值，若是则重复上述修正过程，若否则位置级约束方程修正完毕。

\subsection{速度级违约修正}

广义速度修正量可由下式直接获得，无迭代过程。
\begin{equation}\label{eq:10.69}
  \Delta\dot{\boldsymbol{q}} = -\boldsymbol{c}_q^+ \dot{\boldsymbol{c}}.
\end{equation}
